% ============================================================================
% Vendored preamble for the HyperMIS reproducibility document (main.tex).
%
% Self-contained snapshot of the packages and macros that the generated
% figure/table fragments depend on, distilled from the paper's paper.tex,
% defs.tex and tikz/tikzdefs.tex.  It lets THIS repository build every
% generated artifact into a single PDF WITHOUT the paper checkout.
%
% It is NOT used by the paper (the paper has its own preamble).  If the
% generated fragments start using a new macro, add it here to keep main.tex
% building.  The plot colours are regenerated (write_plot_colors), everything
% else here is static.
% ============================================================================

\usepackage[a4paper, margin=1.5cm]{geometry}
\usepackage[T1]{fontenc}
\usepackage[utf8]{inputenc}
\usepackage{lmodern}

% --- tables -----------------------------------------------------------------
\usepackage{booktabs}
\usepackage{array}
\usepackage{multirow}
\usepackage{multicol}
\usepackage[table]{xcolor}
\usepackage{graphicx}
\usepackage{adjustbox}      % \resizebox around the wide tables
\usepackage{pdflscape}      % landscape per-instance table (tab:overview)

% --- math -------------------------------------------------------------------
\usepackage{amsmath}
\usepackage{amssymb}
\usepackage{amsfonts}
\usepackage{bm}
\usepackage{mathtools}

% --- numbers ----------------------------------------------------------------
\usepackage{siunitx}
\sisetup{detect-weight=true, detect-inline-weight=math,
         group-separator={\,}, group-minimum-digits=4}
\usepackage{numprint}
\npthousandsep{\,}
\npdecimalsign{.}
\usepackage{xstring}        % \IfStrEqCase inside \rname

% --- references (\nameref is the \rname fallback) ---------------------------
\usepackage{hyperref}

% --- plots ------------------------------------------------------------------
\usepackage{pgf}
\usepackage{pgfplots}
\usepackage{pgfplotstable}
\usepackage{tikz}
\usepgfplotslibrary{colorbrewer, statistics, groupplots}
\usetikzlibrary{arrows.meta, positioning, tikzmark, calc, bending,
                shadows, shapes, patterns}
\pgfplotsset{compat=1.18}
\usepackage{wasysym}        % timeout glyph in the performance profile

% Generated plot colours c<name>, refreshed by write_plot_colors().
% Auto-generated plot colours -- do not edit.
% Source: paper_plot_common.write_plot_colors().
\definecolor{cTime}{RGB}{27,158,119}
\definecolor{cnored}{RGB}{27,158,119}
\definecolor{cred}{RGB}{217,95,2}
\definecolor{ctie}{RGB}{117,112,179}
\definecolor{crecompute}{RGB}{27,158,119}
\definecolor{cprecompute}{RGB}{217,95,2}
\definecolor{condemand}{RGB}{117,112,179}
\definecolor{cnrilp}{RGB}{217,95,2}
\definecolor{cfrilp}{RGB}{231,41,138}
\definecolor{crilp}{RGB}{102,166,30}
\definecolor{cilp}{RGB}{230,171,2}
\definecolor{cgfilp}{RGB}{166,118,29}
\definecolor{cgfrilp}{RGB}{102,102,102}
\definecolor{cstruction}{RGB}{27,158,119}
\definecolor{crstruction}{RGB}{217,95,2}
\definecolor{cbmatching}{RGB}{31,120,180}
\definecolor{crbmatching}{RGB}{227,26,28}
\definecolor{cvcsolver}{RGB}{23,190,207}
\definecolor{crvcsolver}{RGB}{152,78,163}
\definecolor{csatreduce}{RGB}{255,127,0}
\definecolor{crsatreduce}{RGB}{65,105,105}
\definecolor{cgrilp}{RGB}{128,0,38}
% Named DARK2 palette shared with the reduction figures / obs boxes.
\definecolor{d2green}{RGB}{27,158,119}
\definecolor{d2orange}{RGB}{217,95,2}
\definecolor{d2purple}{RGB}{117,112,179}
\definecolor{d2pink}{RGB}{231,41,138}
\definecolor{d2lime}{RGB}{102,166,30}
\definecolor{d2gold}{RGB}{230,171,2}
\definecolor{d2brown}{RGB}{166,118,29}
\definecolor{d2gray}{RGB}{102,102,102}


% ============================================================================
% Macros the fragments reference, copied from the paper's defs.tex.
% ============================================================================
\newcommand{\strong}{\textsc{R}}
\newcommand{\strongeduced}{\strong{}educed}
\newcommand{\ilp}{\textsc{ILP}}
\newcommand{\ilpg}{\textsc{ILP$_G$}}
\newcommand{\satAndReduce}{\textsc{SatAndReduce}}
\newcommand{\weGotYouCovered}{\textsc{WeGotYouCovered}}
\newcommand{\struction}{\textsc{Struction}}
\newcommand{\cfg}[1]{\textsc{#1}}
\newcommand{\bmatching}{\cfg{matching}}
\newcommand{\rmode}[2][]{\strong$_{#2}^{#1}$}
\newcommand{\tred}[1]{\ensuremath{t_{\strong_{#1}}^{\mathrm{red}}}}
\newcommand{\ttot}[1]{\ensuremath{t_{\strong_{#1}}}}
\newcommand{\red}[3][]{\strong$_{#2}^{#1}$-#3}
\definecolor{lightgray}{rgb}{0.86,0.86,0.86}   % shaded rows in tab:overview
% configuration name macros (\c... = solver, prefixed r/f/g variants = reduced)
\newcommand{\cilp}{\ilp}
\newcommand{\crilp}{\red{r}{\ilp}}
\newcommand{\cnrilp}{\red{p}{\ilp}}
\newcommand{\cfrilp}{\red{d}{\ilp}}
\newcommand{\cgfilp}{\ilpg}
\newcommand{\cgfrilp}[1][]{\red[\star]{d}{\ilpg}}
\newcommand{\cgrilp}[1][]{\red[\star]{G}{\ilpg}}
\newcommand{\cstruction}{\struction}
\newcommand{\crstruction}[1][]{\red[\star]{d}{\struction}}
\newcommand{\csatreduce}{\satAndReduce}
\newcommand{\crsatreduce}[1][]{\red[\star]{d}{\satAndReduce}}
\newcommand{\cvcsolver}{\weGotYouCovered}
\newcommand{\crvcsolver}[1][]{\red[\star]{d}{\weGotYouCovered}}
\newcommand{\cbmatching}{\bmatching}
\newcommand{\crbmatching}[1][]{\red[T]{d}{\bmatching}}
}

% Revision markup (fragments usually do not use it, kept for completeness).
\definecolor{revred}{rgb}{0.75,0.05,0.05}
\newcommand{\rev}[1]{\textcolor{revred}{#1}}
\newcommand{\revcolor}{\color{revred}}

% Abbreviated rule names used in the tables.  In the paper these link to the
% reduction definitions; standalone there is no Section 4 to link to, so
% \redlink just typesets the short name (no undefined-reference warnings).
\newcommand{\redlink}[2]{#2}
\DeclareRobustCommand{\rname}[1]{%
  \IfStrEqCase{#1}{%
    {edgesizereduction}{\redlink{#1}{Edge Size}}%
    {degree-one}{\redlink{#1}{Degree-One}}%
    {edgedomreduction}{\redlink{#1}{F.\,Edge Dom.}}%
    {fastdomination}{\redlink{#1}{F.\,Vertex Dom.}}%
    {vertexdomination}{\redlink{#1}{Vertex Dom.}}%
    {twinreduction}{\redlink{#1}{Twins}}%
    {unconfinedreduction}{\redlink{#1}{Unconfined}}%
    {edgedominationgeneral}{\redlink{#1}{Edge Dom.}}%
  }[\texttt{#1}]%
}
